\chapter{2020年「大学物理1」\hdu 期末试题}
\date{2020年6月??日}{大学物理教学团队}{A0715011}
{未央物理讲师 \href{https://qm.qq.com/q/UPbGudx8cK}{\textbf{Axia}}}
{\qrcode{https://qm.qq.com/q/asJiytD5pC}}{「未央」HDU 物理营}
{\qrcode{http://weixin.qq.com/r/mhKYgFTEMQlOrRDj90eI}}{未央学社公众号}
\section{选择题（每题3分，共27分）}
{
    \begin{choice}{D}{质点运动学}
        质点做曲线运动，$\boldsymbol{r}$表示位置矢量，$\boldsymbol{v}$表示速度，$\boldsymbol{a}$表示加速度，$S$表示路程，$a_{\tau}$表示切向加速度.下列表达式中
        \begin{multicols}{4}
            \begin{enumerate}[label=(\arabic*)]
                \item $\d v/\d t=a$
                \item $\d r/\d t=v$
                \item $\d S/dt=v$
                \item $\abs{\d\boldsymbol{v}/\d t}=a_{\tau}$
            \end{enumerate}
        \end{multicols}
        \begin{tasks}(4)
            \task 只有(1)、(4)是对的
            \task 只有(2)、(4)是对的
            \task 只有(2)是对的
            \task 只有(3)是对的
        \end{tasks}
    \end{choice}
}
\section{计算题}
\hspace{-2.16em}
\begin{minipage}{0.67\textwidth}
    \begin{exercise}{8}{动量}
        水平小车的B端固定一轻弹簧，弹簧为自然长度时，靠在弹簧上的滑块距小车A端为$l=1.1\mathrm{m}$.已知小车质量$M=10\mathrm{kg}$，滑块质量$m=1\mathrm{kg}$,弹簧的劲度系数$k=110\mathrm{N/m}$.现推动滑块将弹簧压缩$\Delta l=0.05\mathrm{m}$.并维持滑块与小车静止，然后同时释放滑块与小车.忽略一切摩擦.求
        \begin{enumerate}
            \item 滑块与弹簧刚刚分离时,小车及滑块相对地的速度各为多少?
            \item 滑块与弹簧分离后,又经多少时间滑块从小车上掉下来?
        \end{enumerate}
    \end{exercise}
\end{minipage}
\hfill
\begin{minipage}[c]{0.33\textwidth}
\begin{center}
    \begin{tikzpicture}
    
    \end{tikzpicture}
\end{center}
\end{minipage}

\begin{exercise}{5}{角动量守恒}
    质量为$75\mathrm{kg}$的人站在半径为$2\mathrm{m}$的水平转台边缘.转台的固定转轴竖直通过台心且无摩擦.转台绕竖直轴的转动惯量为$3000\mathrm{kg/m^2}$.开始时整个系统静止.现人以相对于地面为$1\mathrm{m/s}$的速率沿转台边缘行走.求人沿转台边缘行走一周，回到他在转台上初始位置所用的时间.
\end{exercise}

\hspace{-2.16em}
\begin{minipage}{0.67\textwidth}

\begin{exercise}{5}{高斯定理}
    真空中一立方体形的高斯面边长$a=0.3\mathrm{m}$，位于图中所示位置.已知空间的场强分布为
    $$E_x=bx,\ E_y=0,\ E_z=0$$
    常量$b=1000\mathrm{N\ab(C\cdot m)}^{-1}$.试求通过该高斯面的电通量.
\end{exercise}
\end{minipage}
\hfill
\begin{minipage}[c]{0.33\textwidth}
\begin{center}
    \begin{tikzpicture}
    \draw [thick,->] (0,0)--(0,2);
    \draw [thick,->] (0,0)--(3.2,0);
    \draw [thick,->] (0,0)--(-1,-1);
    \node [anchor=east] at (0,0) {$O$};
    \node [anchor=north] at (3,0) {$x$};
    \node [anchor=west] at (0,2) {$y$};
    \node [anchor=east] at (-1,-1) {$z$};
    \draw [densely dashed,thick] (1,0)--(2,0)--(2,1)--(1,1)--(1,0)--cycle;
    \draw [densely dashed,thick] (0.646,-0.354)--(1.646,-0.354)--(1.646,0.646)--(0.646,0.646)--(0.646,0.646)--cycle;
    \draw [densely dashed,thick] (0.646,-0.354)--(1,0);
    \draw [densely dashed,thick] (1.646,-0.354)--(2,0);
    \draw [densely dashed,thick] (1.646,0.646)--(2,1);
    \draw [densely dashed,thick] (0.646,0.646)--(1,1);
    \draw [|<->|,thick] (-0.354,-0.554)--(0.646,-0.554) node [midway,below]{$a$};
    \draw [|<->|,thick] (0.646,-0.554)--(1.646,-0.554) node [midway,below]{$a$};
    \draw [|<->|,thick] (1.646,-0.554)--(2.2,0) node [midway,below right]{$a$};
    \draw [|<->|,thick] (2.4,0)--(2.4,1) node [midway,right]{$a$};
    \node [teal,anchor=south] at (1,1) {$\boldsymbol{E}_{x_1}$};
    \node [紺桔梗,anchor=south] at (2,1) {$\boldsymbol{E}_{x_2}$};
    \end{tikzpicture}
\end{center}
\end{minipage}

\begin{exercise}{6}{电场力做功}
    如图所示，一半径为$R$的均匀带正电圆环，其电荷线密度为$\lambda$.在其轴线上有$A$、$B$两点，它们与环心的距离分别为$\overline{OA}=\sqrt{3}R$，$\overline{OB}=\sqrt{8}R$.一质量为$m$、电荷为$q$的粒子从$A$点运动到$B$点.求在此过程中电场力所做的功.
\end{exercise}

\hspace{-2.16em}
\begin{minipage}{0.67\textwidth}
    \begin{exercise}{6}{磁感应强度}
        如图所示，半径为$R$，$\lambda(>0)$的均匀带电的圆线圈绕过圆心与圆平面的轴以角速度$\omega$转动.求轴线上任一点的$\boldsymbol{B}$的大小及其方向.
    \end{exercise}
\end{minipage}
\hfill
\begin{minipage}[c]{0.33\textwidth}
\begin{center}
    \begin{tikzpicture}
    
    \end{tikzpicture}
\end{center}
\end{minipage}

\hspace{-2.16em}
\begin{minipage}{0.67\textwidth}
    \begin{exercise}{8}{电磁感应}
        如图所示，一长直导线中通有电流$I$，有一垂直于导线、长度为$l$的金属棒$AB$在包含导线的平面内，以恒定的速度$v$沿与棒成$\theta$角的方向移动.开始时，棒的$A$端到导线的距离为$d$，求任意时刻金属棒中的动生电动势，并指出棒哪端的电势高.
    \end{exercise}
\end{minipage}
\hfill
\begin{minipage}[c]{0.33\textwidth}
\begin{center}
    \begin{tikzpicture}
    
    \end{tikzpicture}
\end{center}
\end{minipage}